\documentclass[11pt,a4paper]{article}

\usepackage[utf8]{inputenc}
\usepackage[T1]{fontenc}
\usepackage{lmodern}
\usepackage[margin=2.0cm]{geometry}
\usepackage{amsmath}
\usepackage{booktabs}
\usepackage{array}
\usepackage{tabularx}
\usepackage{graphicx}
\usepackage{xcolor}
\usepackage{hyperref}
\usepackage{microtype}
\usepackage{caption}
\usepackage{float}

\definecolor{navy}{HTML}{102A43}
\definecolor{blue}{HTML}{1677B8}
\hypersetup{
  colorlinks=true,
  linkcolor=blue,
  citecolor=blue,
  urlcolor=blue,
  pdftitle={From Transactions to Buildings: Leakage-Aware External Data Enrichment for Paris Apartment Valuation},
  pdfauthor={Badr Ettousy}
}
\graphicspath{{../../visuals/phase3/}}
\captionsetup{font=small,labelfont=bf}
\setlength{\parskip}{0.35em}
\setlength{\parindent}{0pt}
\renewcommand{\arraystretch}{1.14}

\title{\textcolor{navy}{\textbf{From Transactions to Buildings}}\\[0.25em]
\large Leakage-Aware External Data Enrichment for\\
Paris Apartment Valuation}
\author{Badr Ettousy}
\date{August 2026}

\begin{document}
\maketitle

\begin{abstract}
This study evaluates whether authoritative building and neighborhood data can
improve an automated valuation model (AVM) previously limited to French
Demandes de valeurs fonci\`eres (DVF) transactions and lagged comparable sales.
The fixed cohort contains 143,009 ordinary Paris apartment sales from
2021--2025. Phase 3 links each transaction to Base Adresse Nationale (BAN)
address identity; Base de Donn\'ees Nationale des B\^atiments (BDNB) building
attributes and dated energy-performance diagnostics (DPE); rail stations;
schools, parks and nearby services; and strategic transport-noise maps. Exact
BAN matches cover 99.49\% of sales, 99.997\% receive a BDNB building identity,
and 76.85\% receive a DPE that existed by the transaction date. No future-DPE
join is observed. Configuration selection uses 2024 only; selected models are
refitted on 2021--2024 and evaluated on 28,330 sales from 2025. The selected
CatBoost model obtains EUR 93,136 mean absolute error (MAE), EUR 192,806 root
mean squared error, $R^2=0.8679$, and 11.70\% median absolute percentage error.
Relative to the frozen Phase 2 comparable-sale correction (EUR 96,428 MAE), the
MAE reduction is EUR 3,293, or 3.42\%, with a paired 95\% bootstrap interval of
EUR 2,735--3,846. Dated DPE, building and neighborhood ablations each improve
MAE, although their gains are not additive or causal. Current BAN/BDNB and
amenity snapshots are explicitly classified as retrospective reconstruction;
only dated DPE and lagged comparable-sale variables satisfy full point-in-time
rules. The result shows a meaningful but bounded improvement and does not turn
the AVM into a certified appraisal.
\end{abstract}

\textbf{Keywords:} automated valuation model; DVF; BAN; BDNB; DPE; entity
resolution; point-in-time features; CatBoost; Paris real estate.

\section{Executive result}

Phase 3 was designed to answer one practical question: does adding reliable
property identity, building condition proxies, energy information, access and
neighborhood context improve out-of-year valuation accuracy beyond Phase 2?
The answer is yes on the fixed 2025 cohort. The selected model reduces MAE by
EUR 3,293 (3.42\%), improves MdAPE from 12.22\% to 11.70\%, increases the share
of estimates within 20\% from 71.69\% to 73.39\%, and raises $R^2$ from 0.8601
to 0.8679. The paired interval for the MAE reduction excludes zero.

The contribution is broader than one score. It includes a versioned Bronze
source layer, deterministic address and building resolution, a row-level match
audit, a leakage-aware DPE join, an enriched 141-column Gold table, controlled
feature-family ablations, a reusable CatBoost artifact and an auditable
single-property inference path.

\begin{figure}[H]
\centering
\includegraphics[width=\textwidth]{phase3_results_dashboard.png}
\caption{Executive Phase 3 result. Error bars are nonparametric 95\% bootstrap
intervals. Paired improvement intervals resample transaction-level differences
against the same frozen Phase 2 predictions.}
\label{fig:executive}
\end{figure}

\section{Problem setting and prior phases}

The target is total transaction value in euros for an ordinary apartment sale
in Paris. Phase 1 established a chronological tabular-model benchmark using
DVF characteristics. Phase 2 added Lambert-93 spatial variables, H3 cells and
historical comparable-sale statistics under a minimum 90-day information lag.
Its best method combined 75\% of a LightGBM estimate with 25\% of the trailing
365-day, 500-metre local median price per square metre. That correction reached
EUR 96,428 MAE in 2025, improving the Phase 1 reference only modestly.

Phase 2 also exposed the main information gap. DVF describes the transaction
and apartment size but not the building, dated energy state, exact persistent
address identity, transit access or environmental context. Phase 3 therefore
focuses on external data linkage rather than a more exotic neural architecture.

\section{Data and external sources}

\subsection{Fixed transaction cohort}

The transaction cohort is unchanged from the earlier phases. DVF rows are
collapsed and filtered at mutation level so that the target is not duplicated
across property lines \cite{dvf}. The cohort retains arm's-length sales containing exactly
one apartment, while allowing dependencies such as cellars or parking spaces.
Required fields include date, value, surface and Paris coordinates. Fixed
plausibility bounds are 9--300 m$^2$ for surface, EUR 50,000--10,000,000 for
value and EUR 2,000--30,000 per m$^2$.

\begin{table}[H]
\centering
\caption{Chronological cohorts used throughout Phase 3.}
\label{tab:splits}
\begin{tabular}{lrrl}
\toprule
Period & Sales & Cumulative training sales & Role \\
\midrule
2021 & 30,397 & 30,397 & Development \\
2022 & 33,080 & 63,477 & Development \\
2023 & 26,798 & 90,275 & Development \\
2024 & 24,404 & 114,679 & Configuration selection \\
2025 & 28,330 & --- & Frozen chronological test \\
\midrule
Total & 143,009 & --- & Fixed Gold identity \\
\bottomrule
\end{tabular}
\end{table}

\subsection{Authoritative enrichment sources}

Table~\ref{tab:sources} summarizes the new inputs. Every downloaded artifact is
stored under \path{data/bronze/} with producer, licence, source URL, retrieval
timestamp, byte size and SHA-256 recorded in
\path{reports/phase3/source_manifest.json}. The BDNB departmental archive is
version 2026-02.a for Paris and contains both DPE records and the official
DPE-to-building relation; a second bulk ADEME download is therefore not needed.

\begin{table}[H]
\centering
\caption{Phase 3 external sources and snapshot interpretation.}
\label{tab:sources}
\footnotesize
\begin{tabularx}{\textwidth}{p{3.05cm}X p{3.25cm}}
\toprule
Source / producer / licence & Variables introduced & Temporal status \\
\midrule
BAN \cite{ban}\\IGN, Licence Ouverte 2.0 & Normalized address ID, FANTOIR road identity, coordinates and cadastral references & Current static reference \\
BDNB 2026-02.a \cite{bdnb}\\CSTB, Licence Ouverte 2.0 & Building ID, footprint, levels, height, year/period, dwelling and lot counts, materials and coarse state & Retrospective static reconstruction \\
DPE in BDNB \cite{dpe}\\ADEME/CSTB, Licence Ouverte 2.0 & Energy and GHG classes, consumption, emissions, heating, surface and diagnostic date & Point-in-time date filter \\
Rail stations \cite{idfm}\\\^Ile-de-France Mobilit\'es, Licence Ouverte 2.0 & Metro, RER, train and tram distance/density by mode and line & Current static context \\
Schools and parks \cite{paris}\\Ville de Paris, ODbL & Nearest distance, counts and green-space area & Current static context \\
Services\\BDNB BPE, Licence Ouverte 2.0 & Nearby shop/service distance and counts & Retrospective static context \\
Noise \cite{noise}\\DRIEAT/CEREMA, Licence Ouverte 2.0 & 2022 road, RATP-rail and SNCF-rail Lden bands & Fixed 2022 strategic map \\
\bottomrule
\end{tabularx}
\end{table}

\section{Data architecture and entity resolution}

The pipeline follows Bronze--Silver--Gold layering (Figure~\ref{fig:pipeline}).
Bronze preserves source snapshots. Silver resolves identities and records how
every match was obtained. Gold retains one row per original DVF sale and adds
the external fields plus explicit missingness and provenance indicators.

\begin{figure}[H]
\centering
\includegraphics[width=\textwidth]{phase3_data_pipeline.png}
\caption{Source-to-model architecture. The feature contract distinguishes
dated information from retrospective static reconstruction.}
\label{fig:pipeline}
\end{figure}

\subsection{Exact address identity}

The primary address key combines commune, the four-character FANTOIR/DVF road
code, street number and normalized suffix. It is deterministic and avoids
fuzzy street-name matching. The key resolves to the BAN interoperability ID.
The resulting exact BAN match rate is 99.49\%. Coordinates and labels from BAN
are retained for audit, not used to overwrite the original transaction identity
silently.

\subsection{BAN-to-BDNB building resolution}

The official BAN--BDNB relationship table can return several buildings for one
address. Candidates are ranked using the relationship reliability field; the
highest-ranked building is selected deterministically. When the address route
is unavailable, the cadastral parcel relation is used as a fallback. The Gold
table records the selected BDNB building, match method, confidence, number of
candidates and both address and parcel evidence. This produces a 99.997\%
building match rate and 99.46\% high-confidence rate. The row-level Silver audit
is stored in \path{data/silver/phase3_match_audit.parquet}.

\subsection{Identity-preserving joins}

The build asserts that the enriched table has exactly 143,009 rows, that the
original mutation order and IDs are preserved and that no duplicate mutation
is created. This matters because one-to-many building, DPE or amenity joins can
otherwise inflate the training sample and create false accuracy. The delivered
Gold table has 141 typed columns and a recorded SHA-256 contract.

\section{Feature engineering and temporal contract}

\subsection{Building structure, age and condition}

BDNB contributes construction year or period, age at sale, level count, height,
footprint, dwelling count, total/parking lot counts, wall and roof materials,
usage class and coarse building state. Construction age is computed relative
to each sale year and invalid negative or implausible values are rejected. The
median building age in the 2025 test is 125 years, reflecting the old Paris
housing stock. These fields are valuable retrospective descriptors, but the
2026 snapshot does not prove that every value was observable in earlier years.

\subsection{Point-in-time DPE}

For a sale at date $t$, a diagnostic is eligible only when
\begin{equation}
  \mathrm{date\_etablissement\_dpe} \leq t.
\end{equation}
Within each resolved building, an as-of join selects the latest eligible DPE.
The pipeline records DPE age in days, apartment-surface discrepancy and match
confidence. A hard validation rejects any future diagnostic. The observed
future-DPE leakage count is zero.

Overall historical coverage is 76.85\%, but it is strongly time-dependent:
13.12\% in 2021, 83.38\% in 2022, 97.10\% in 2023, 99.05\% in 2024 and 99.33\%
in 2025. This pattern is expected because the current DPE regime begins in the
study period; it also explains why a single aggregate coverage percentage is
insufficient.

\subsection{Accessibility, amenities and noise}

Coordinates are transformed to Lambert-93 before metric distance calculations.
The new context variables include nearest Metro, RER, train, tram, school,
green-space centroid and shop/service distances; counts of stations, modes,
lines, schools, green spaces and services inside fixed radii; green-space area;
and maximum intersecting road/RATP/SNCF Lden bands. In the 2025 cohort, median
straight-line distances are 222 m to Metro, 946 m to RER, 154 m to a school,
103 m to a green-space centroid and 6 m to a BDNB BPE service location.

Noise values represent strategic contour bands rather than parcel-level acoustic
measurements. Locations outside a mapped high-noise contour are represented by
the low/no-band category and a missing-contour indicator. They must not be read
as an exact measured decibel value.

\begin{figure}[H]
\centering
\includegraphics[width=\textwidth]{phase3_enrichment_profile.png}
\caption{External-feature profile on the 2025 test cohort. DPE shares are
conditional on a diagnostic available by the sale date. Distances are
straight-line Lambert-93 measures. Noise is a strategic contour category.}
\label{fig:profile}
\end{figure}

\subsection{Unavailable requested fields}

The reviewed public BAN, BDNB and DPE fields do not provide a reliable
apartment-floor or elevator variable for this cohort. Building level count is
not a substitute for apartment floor. Gold therefore includes explicit
\texttt{apartment\_floor\_available=0} and
\texttt{elevator\_source\_available=0} flags while leaving the values missing.
This avoids converting an unsupported assumption into model input. A lawful
notarial, copropri\'et\'e, property-manager or licensed-listing source would be
needed to complete those fields.

\subsection{Information-time classification}

The final contract separates two evidence classes:
\begin{itemize}
  \item \textbf{Point-in-time:} DPE dated on or before the sale and all Phase 2
  comparable-sale features subject to a minimum 90-day publication lag.
  \item \textbf{Retrospective static:} current BAN identity, most BDNB physical
  attributes, current transport/school/park/service snapshots and the fixed
  2022 strategic noise map.
\end{itemize}
The full model is consequently a strong retrospective reconstruction, not a
claim that every external value was operationally available on each historical
valuation date.

\section{Modeling and evaluation}

\subsection{Feature-family ablations}

All candidates retain the core DVF and Phase 2 variables so that Phase 3 tests
incremental information. Five LightGBM feature sets are compared: Phase 2 only,
dated DPE, building/linkage, neighborhood context and the full enrichment
without raw high-cardinality identity. Each is evaluated with MAE and Huber
objectives using a small predefined configuration grid. CatBoost additionally
receives raw BAN address and BDNB building categories because its ordered
categorical handling is suitable for high-cardinality identity. The final
CatBoost configuration has 650 iterations, learning rate 0.05, depth 7,
$\ell_2$ leaf regularization 5 and MAE loss.

The CatBoost comparison changes both model family and access to raw identity.
Its advantage over full LightGBM therefore cannot be attributed to identity
alone. Likewise, feature-family ablation gains overlap and must not be added.

\subsection{Chronological protocol}

Development uses 2021--2023. The 2024 cohort selects the feature set, objective
and hyperparameters. Selected candidates are refitted once on 2021--2024 and
scored on the common 2025 rows. No 2025 record is used by a fitting or selection
call. However, the 2025 cohort had already been examined in Phases 1 and 2; it is
frozen for Phase 3 computation but not a pristine analyst-blind external test.
A later prospective cohort is required for final deployment evidence.

\subsection{Metrics and uncertainty}

The primary metric is
\begin{equation}
  \mathrm{MAE}=\frac{1}{n}\sum_{i=1}^{n}|y_i-\hat y_i|.
\end{equation}
Secondary metrics are RMSE, $R^2$, mean and median absolute percentage error
(MAPE and MdAPE), and the shares within 10\% and 20\% of the transaction value.
MAE intervals use nonparametric bootstrap resampling. Improvement intervals use
paired resampling of per-sale absolute-error differences, preserving the fact
that every candidate predicts the same transactions.

\section{Results}

\subsection{Ablation and final test performance}

Table~\ref{tab:results} reports the principal 2025 results. Every MAE-objective
enrichment ablation improves on the frozen Phase 2 reference. Dated DPE alone
reduces MAE by EUR 1,206; building/linkage by EUR 1,858; neighborhood context by
EUR 796; and the full no-identity LightGBM model by EUR 2,460. Their paired 95\%
intervals exclude zero. The selected CatBoost model improves MAE by EUR 3,293.

\begin{table}[H]
\centering
\caption{Frozen 2025 chronological test. Intervals are paired MAE reductions
against Phase 2; positive values favor the row method.}
\label{tab:results}
\small
\begin{tabular}{lrrrr}
\toprule
Method & MAE (EUR) & MdAPE & Within 20\% & Reduction (95\% CI), EUR \\
\midrule
CatBoost, full identity/enrichment & \textbf{93,136} & \textbf{11.70\%} & \textbf{73.39\%} & \textbf{3,293 (2,735--3,846)} \\
LightGBM, full without raw identity & 93,968 & 11.80\% & 73.03\% & 2,460 (1,860--3,137) \\
LightGBM, building/linkage & 94,571 & 11.90\% & 72.28\% & 1,858 (1,271--2,493) \\
LightGBM, dated DPE & 95,222 & 11.97\% & 72.28\% & 1,206 (657--1,867) \\
LightGBM, neighborhood context & 95,632 & 11.99\% & 71.74\% & 796 (293--1,355) \\
Phase 2 comparable correction & 96,428 & 12.22\% & 71.69\% & Reference \\
\bottomrule
\end{tabular}
\end{table}

For the winner, RMSE is EUR 192,806, $R^2=0.8679$, MAPE is 18.12\%, 43.68\%
of estimates are within 10\% and 73.39\% are within 20\%. Phase 3 reduces
absolute error on 54.21\% of individual test sales; the median transaction-level
gain is EUR 1,496. The improvement therefore reflects many modest gains plus a
smaller number of large corrections, not universal superiority on every sale.

\subsection{Calibration, value segments and time stability}

Figure~\ref{fig:diagnostics} shows that the predicted and observed values remain
tightly aligned over most of the range, with dispersion increasing among
high-value sales. Phase 3 lowers MAE in every transaction-value decile shown,
including the upper decile where absolute errors are much larger. Monthly 2025
MAE also remains below Phase 2 in every month. This supports temporal stability
within the test year, although it does not measure a new market regime.

\begin{figure}[H]
\centering
\includegraphics[width=\textwidth]{phase3_prediction_diagnostics.png}
\caption{Frozen-test diagnostics. The transaction-level gain distribution is
clipped for legibility; reported averages use the unclipped differences.}
\label{fig:diagnostics}
\end{figure}

\subsection{Geographic heterogeneity}

The mean improvement is positive in 19 of 20 arrondissements. It is largest in
the 6th, 4th, 7th, 8th and 1st arrondissements and approximately neutral in the
16th (EUR $-167$ mean change). The map also shows intermixed gains and losses
within the same neighborhoods. This is consistent with heterogeneous errors at
the individual property level and argues against interpreting arrondissement
averages as causal amenity premiums.

\begin{figure}[H]
\centering
\includegraphics[width=\textwidth]{phase3_spatial_performance.png}
\caption{Spatial distribution of transaction-level changes and mean MAE
reduction by arrondissement. Green marks lower absolute error under Phase 3.}
\label{fig:spatial}
\end{figure}

\subsection{Model use of enriched information}

CatBoost prediction-value-change importance assigns 79.68\% of total importance
to core DVF/spatial variables, dominated by surface and log surface. Historical
comparables account for 12.95\%; building/linkage 2.30\%; neighborhood context
2.01\%; DPE/energy 1.86\%; and identity/micro-location 1.20\%. The strong
ablation gains alongside small individual importances are not contradictory:
external variables can refine difficult cases while surface still explains
most price-level variation. Importance values describe the fitted model's use
of features and do not estimate willingness to pay or causal effects.

\begin{figure}[H]
\centering
\includegraphics[width=\textwidth]{phase3_feature_importance.png}
\caption{CatBoost prediction-value-change importance for the selected model.
Feature-family totals sum to 100. Correlated variables may divide importance.}
\label{fig:importance}
\end{figure}

\section{Interpretation and practical value}

Three conclusions follow. First, external information addresses a real omitted-
variable problem. Each enrichment family provides an independently positive
test signal under the controlled LightGBM ablations. Second, reliable identity
is operationally important even when its direct feature importance is small: it
allows building, DPE and context records to attach to the correct transaction
and enables CatBoost to learn repeated micro-location patterns. Third, the
improvement is meaningful but bounded. A 3.42\% MAE reduction does not remove
unobserved apartment condition, renovation quality, floor, elevator, view,
exposure, occupancy or private outdoor space.

The selected artifact supports single-property inference through
\path{paris_avm.inference.phase3}. The caller supplies the apartment attributes,
coordinates, valuation date, commune code and FANTOIR street code. The command
resolves an exact existing BAN/BDNB identity, rolls DPE back to the valuation
date, recomputes lagged comparable features and returns an automated indication.
It also warns that building and neighborhood fields use the 2026 static snapshot.

\section{Limitations and threats to validity}

\begin{enumerate}
  \item \textbf{Static-snapshot bias.} Current physical/context snapshots can
  improve reconstruction of past sales without proving historical availability.
  This is the most important limitation.
  \item \textbf{Previously inspected test year.} The 2025 rows are excluded
  from Phase 3 fitting and selection, but earlier phases had exposed aggregate
  2025 performance. Prospective confirmation is still necessary.
  \item \textbf{Residual property heterogeneity.} Floor, elevator, interior
  condition, view and renovation quality remain unavailable.
  \item \textbf{DPE linkage.} A building can contain several apartments and
  diagnostics. The latest eligible building-linked DPE may not be the exact
  apartment; surface discrepancy and confidence fields make this uncertainty
  visible but do not eliminate it.
  \item \textbf{Noise and amenity semantics.} Straight-line distances ignore
  walking networks and barriers. Strategic noise contours are categorical
  planning layers, not indoor measurements.
  \item \textbf{Bootstrap scope.} The paired interval measures sampling
  variation within the 2025 cohort, not uncertainty under future interest-rate,
  policy or market-regime shifts.
  \item \textbf{Non-causal analysis.} Ablations and feature importance quantify
  predictive association only.
\end{enumerate}

\section{Reproducibility and transfer package}

The complete Phase 3 transfer package is contained in the project:
\begin{itemize}
  \item acquisition: \path{paris_avm.data.acquire_phase3};
  \item entity resolution and features: \path{paris_avm.features.phase3};
  \item ablation/model benchmark: \path{paris_avm.modeling.benchmark_phase3};
  \item figures: \path{paris_avm.visualization.phase3};
  \item inference: \path{paris_avm.inference.phase3};
  \item Gold table: \path{data/gold/phase3_sale_features.parquet};
  \item audit and contracts: \path{data/silver/phase3_match_audit.parquet} and
  \path{reports/phase3/};
  \item selected model: \path{models/phase3/phase3_selected_model.joblib}.
\end{itemize}

The deterministic execution order is:
\begin{verbatim}
python -m paris_avm.data.acquire_phase3 --skip-dpe
python -m paris_avm.features.phase3
python -m paris_avm.modeling.benchmark_phase3
python -m paris_avm.visualization.phase3
powershell -ExecutionPolicy Bypass -File docs/paper/compile_phase3.ps1
\end{verbatim}
The \texttt{--skip-dpe} option is intentional because the dated ADEME DPE rows
and official DPE-to-building relation are already present in the versioned BDNB
archive.

\section{Conclusion}

Phase 3 converts a transaction-only benchmark into an auditable building-aware
valuation pipeline. Near-complete address/building resolution, strict dated-DPE
selection and controlled feature-family experiments produce a statistically
consistent EUR 3,293 MAE reduction on 28,330 sales from 2025. The strongest
model remains driven by surface and historical market evidence, while external
building, energy and context features provide useful refinements. The proper
next step is not more retrospective tuning: it is prospective evaluation on a
new cohort with vintage-correct snapshots, followed by lawful acquisition of
apartment floor, elevator and condition data. Until then, the output should be
used as a transparent automated indication rather than a certified appraisal.

\begin{thebibliography}{9}

\bibitem{dvf}
Direction g\'en\'erale des Finances publiques.
\emph{Demandes de valeurs fonci\`eres (DVF)}.
\url{https://www.data.gouv.fr/fr/datasets/demandes-de-valeurs-foncieres/}.

\bibitem{ban}
Institut national de l'information g\'eographique et foresti\`ere.
\emph{Base Adresse Nationale: t\'el\'echargements}.
\url{https://adresse.data.gouv.fr/outils/telechargements}.

\bibitem{bdnb}
Centre scientifique et technique du b\^atiment.
\emph{Base de Donn\'ees Nationale des B\^atiments}.
\url{https://www.data.gouv.fr/fr/datasets/base-de-donnees-nationale-des-batiments/}.

\bibitem{dpe}
ADEME.
\emph{Diagnostics de performance \'energ\'etique pour les logements}.
\url{https://data.ademe.fr/datasets/meg-83tjwtg8dyz4vv7h1dqe}.

\bibitem{idfm}
\^Ile-de-France Mobilit\'es.
\emph{Gares et stations du r\'eseau ferr\'e d'\^Ile-de-France par ligne}.
\url{https://data.iledefrance.fr/explore/dataset/gares-et-stations-du-reseau-ferre-dile-de-france-par-ligne/}.

\bibitem{paris}
Ville de Paris.
\emph{Paris Data: \'etablissements scolaires et espaces verts}.
\url{https://opendata.paris.fr/}.

\bibitem{noise}
DRIEAT \^Ile-de-France.
\emph{Cartes de bruit strat\'egiques pour Paris}.
\url{https://www.data.gouv.fr/fr/datasets/lot-de-donnees-relatives-aux-cartes-de-bruit-strategiques-pour-paris/}.

\bibitem{catboost}
L. Prokhorenkova, G. Gusev, A. Vorobev, A. Dorogush, and A. Gulin.
CatBoost: unbiased boosting with categorical features.
\emph{Advances in Neural Information Processing Systems}, 31, 2018.

\bibitem{lightgbm}
G. Ke, Q. Meng, T. Finley, T. Wang, W. Chen, W. Ma, Q. Ye, and T.-Y. Liu.
LightGBM: a highly efficient gradient boosting decision tree.
\emph{Advances in Neural Information Processing Systems}, 30, 2017.

\end{thebibliography}

\end{document}
